\chapter*{Abstract}
\addcontentsline{toc}{chapter}{Abstract}

Il presente elaborato descrive la progettazione e lo sviluppo di \textit{JarFin}, un sistema software per la gestione e l’analisi delle finanze personali dotato di interfaccia conversazionale. L’obiettivo del progetto è definire e realizzare un’architettura software modulare che integri servizi di gestione dei dati finanziari, componenti di analisi algoritmica e un’interfaccia basata su linguaggio naturale, garantendo coerenza, scalabilità e chiarezza nella separazione delle responsabilità tra i diversi moduli.

Il sistema è stato sviluppato seguendo una metodologia Agile, con particolare riferimento a SCRUM e all’approccio Agile Model Driven Development (AMDD), consentendo una progettazione incrementale basata su modelli e una validazione progressiva dei requisiti. L’architettura risultante è composta da servizi distinti, tra cui un servizio di accounting responsabile della gestione delle transazioni finanziarie e un servizio di analytics dedicato all’elaborazione, all’aggregazione e all’analisi dei dati, con particolare attenzione alla progettazione degli algoritmi e alla loro complessità computazionale.

Un elemento caratterizzante del sistema è l’integrazione di un’interfaccia conversazionale che consente all’utente di interrogare il sistema mediante linguaggio naturale. Questa componente introduce specifiche problematiche architetturali, tra cui la gestione dei flussi informativi tra servizi, la traduzione delle richieste in operazioni computazionali e la produzione di risposte coerenti e semanticamente corrette.

Il lavoro presenta inoltre una valutazione dell’architettura progettata e delle soluzioni implementate, evidenziandone i punti di forza e i limiti. JarFin rappresenta un caso di studio concreto per l’applicazione dei principi di progettazione algoritmica, modellazione del software e progettazione architetturale nella realizzazione di un sistema software complesso e modulare.